%% ─────────────────────────────────────────────────────────────────────────────
%% CLASE emate-ucr  —  examen parcial
%% ─────────────────────────────────────────────────────────────────────────────
%%
%% Ejemplo completo con todos los comandos y entornos disponibles.
%% Ver también ejemplo_ejercicios.tex y ejemplo_prueba_corta.tex.
%%
%% COMPILACIÓN
%%   Requiere pdflatex (no xelatex ni lualatex):
%%     pdflatex -synctex=1 -interaction=nonstopmode ejemplo_examen.tex
%%
%%   Para generar la versión con soluciones, crear un archivo aparte:
%%     \PassOptionsToClass{soluciones}{emate-ucr}
%%     %% ─────────────────────────────────────────────────────────────────────────────
%% CLASE emate-ucr  —  examen parcial
%% ─────────────────────────────────────────────────────────────────────────────
%%
%% Ejemplo completo con todos los comandos y entornos disponibles.
%% Ver también ejemplo_ejercicios.tex y ejemplo_prueba_corta.tex.
%%
%% COMPILACIÓN
%%   Requiere pdflatex (no xelatex ni lualatex):
%%     pdflatex -synctex=1 -interaction=nonstopmode ejemplo_examen.tex
%%
%%   Para generar la versión con soluciones, crear un archivo aparte:
%%     \PassOptionsToClass{soluciones}{emate-ucr}
%%     %% ─────────────────────────────────────────────────────────────────────────────
%% CLASE emate-ucr  —  examen parcial
%% ─────────────────────────────────────────────────────────────────────────────
%%
%% Ejemplo completo con todos los comandos y entornos disponibles.
%% Ver también ejemplo_ejercicios.tex y ejemplo_prueba_corta.tex.
%%
%% COMPILACIÓN
%%   Requiere pdflatex (no xelatex ni lualatex):
%%     pdflatex -synctex=1 -interaction=nonstopmode ejemplo_examen.tex
%%
%%   Para generar la versión con soluciones, crear un archivo aparte:
%%     \PassOptionsToClass{soluciones}{emate-ucr}
%%     %% ─────────────────────────────────────────────────────────────────────────────
%% CLASE emate-ucr  —  examen parcial
%% ─────────────────────────────────────────────────────────────────────────────
%%
%% Ejemplo completo con todos los comandos y entornos disponibles.
%% Ver también ejemplo_ejercicios.tex y ejemplo_prueba_corta.tex.
%%
%% COMPILACIÓN
%%   Requiere pdflatex (no xelatex ni lualatex):
%%     pdflatex -synctex=1 -interaction=nonstopmode ejemplo_examen.tex
%%
%%   Para generar la versión con soluciones, crear un archivo aparte:
%%     \PassOptionsToClass{soluciones}{emate-ucr}
%%     \input{ejemplo_examen}
%%
\documentclass{emate-ucr}
% \documentclass[soluciones]{emate-ucr}   % <-- versión con soluciones

%% ── Metadatos obligatorios ───────────────────────────────────────────────────
%%
%% \curso{texto}    Centro del encabezado de página. \\ para saltos de línea.
%% \encabezado{texto}  Título en negrita al inicio del documento.
%%
\curso{I Ciclo 2026\\MA-1022\\Cálculo para Ciencias Económicas II}
\encabezado{Primer Examen Parcial}

%% ── Metadatos opcionales ─────────────────────────────────────────────────────
%%
%% Cualquier subconjunto de estos tres puede usarse o no.
%% Los que estén definidos aparecen en una línea bajo el título:
%%   Fecha: ...  |  Duración: ...  |  Valor: ...
%%
\fecha{15 de abril de 2026}
\duracion{2 horas}
\valor{100 puntos}

\begin{document}

%% \imprimirtitulo
%%   Imprime el encabezado del documento: título y línea de metadatos
%%   (si alguno está definido). Debe ir inmediatamente después de
%%   \begin{document}.
%%
\imprimirtitulo

%% \datosestudiante
%%   Línea para nombre y carné del estudiante. Omitir en hojas de
%%   ejercicios que no requieren identificación individual.
%%
\datosestudiante

%% ── Entorno instrucciones ─────────────────────────────────────────────────────
%%
%% Caja enmarcada con lista numerada de indicaciones.
%% Cada indicación va con \item.
%%
%% Alternativa más simple (sin caja, texto libre):
%%   \begin{indicaciones}
%%     Texto de las indicaciones...
%%   \end{indicaciones}
%%
\begin{instrucciones}
  \item Este examen es individual, con libro y apuntes cerrados.
  \item Justifique cada respuesta. Una respuesta sin desarrollo no se
        calificará aunque sea correcta.
  \item Está prohibido el uso de calculadora.
  \item Si necesita más espacio, use el reverso de la hoja e indíquelo
        claramente.
  \item La deshonestidad académica se reportará según el Reglamento
        Estudiantil de la Universidad de Costa Rica.
\end{instrucciones}

%% ── Entorno ejercicio ────────────────────────────────────────────────────────
%%
%% \begin{ejercicio}[N]   N es el valor en puntos (opcional).
%% Numeración automática. Puede combinarse con \begin{subejercicios}
%% para incisos a), b), c), ...
%%
\begin{ejercicio}[25]
  Resuelva el siguiente sistema de ecuaciones aplicando eliminación de
  Gauss-Jordan. Dé el conjunto solución.
  \[
    \begin{cases}
       x +  y +  z =  6 \\
      2x + 3y +  z = 11 \\
       x + 2y + 3z = 14
    \end{cases}
  \]

  %% \begin{solucion}...\end{solucion}
  %%   Visible solo con la opción [soluciones]. Sin ella, no genera output.
  %%
  \begin{solucion}
    La matriz aumentada es
    \[
      \left(\begin{array}{ccc|r} 1 & 1 & 1 & 6 \\ 2 & 3 & 1 & 11 \\ 1 & 2 & 3 & 14
      \end{array}\right)
      \longrightarrow \cdots \longrightarrow
      \left(\begin{array}{ccc|r} 1 & 0 & 0 & 1 \\ 0 & 1 & 0 & 2 \\ 0 & 0 & 1 & 3
      \end{array}\right).
    \]
    El conjunto solución es $S = \{(1,\,2,\,3)\}$.
  \end{solucion}
\end{ejercicio}

\begin{ejercicio}[25]
  Determine el rango de cada matriz.

  \begin{subejercicios}
    \item $\begin{pmatrix} 1 & 2 & 3 \\ 2 & 5 & 7 \\ 1 & 3 & 4 \end{pmatrix}$

    \begin{solucion}
      El rango es $\mathbf{2}$.
    \end{solucion}

    \item $\begin{pmatrix} 1 & 3 & 2 \\ 2 & 6 & 4 \end{pmatrix}$

    \begin{solucion}
      El rango es $\mathbf{1}$.
    \end{solucion}
  \end{subejercicios}
\end{ejercicio}

\begin{ejercicio}[25]
  Determine si el siguiente sistema homogéneo tiene solución no trivial.
  Justifique.
  \[
    \begin{cases}
       x +  y +  z = 0 \\
      2x + 3y + 4z = 0 \\
       x + 2y + 3z = 0
    \end{cases}
  \]

  \begin{solucion}
    Reduciendo: el sistema tiene solución no trivial (variable libre).
    El conjunto solución es $S = \{(t,\,-2t,\,t) : t \in \mathbb{R}\}$.
  \end{solucion}
\end{ejercicio}

\begin{ejercicio}[25]
  Demuestre que si $A$ es una matriz de $3\times 3$ con rango~$3$, entonces el
  sistema $A\mathbf{x} = \mathbf{b}$ tiene solución única para todo
  $\mathbf{b} \in \mathbb{R}^3$.

  \begin{solucion}
    Si $\operatorname{rango}(A) = 3$, la forma escalonada reducida de $A$ es la
    identidad $I_3$. Para cualquier $\mathbf{b}$, la eliminación de Gauss-Jordan
    produce una única solución $\mathbf{x} = A^{-1}\mathbf{b}$.
  \end{solucion}
\end{ejercicio}

\end{document}

%%% Local Variables:
%%% mode: LaTeX
%%% TeX-master: t
%%% End:

%%
\documentclass{emate-ucr}
% \documentclass[soluciones]{emate-ucr}   % <-- versión con soluciones

%% ── Metadatos obligatorios ───────────────────────────────────────────────────
%%
%% \curso{texto}    Centro del encabezado de página. \\ para saltos de línea.
%% \encabezado{texto}  Título en negrita al inicio del documento.
%%
\curso{I Ciclo 2026\\MA-1022\\Cálculo para Ciencias Económicas II}
\encabezado{Primer Examen Parcial}

%% ── Metadatos opcionales ─────────────────────────────────────────────────────
%%
%% Cualquier subconjunto de estos tres puede usarse o no.
%% Los que estén definidos aparecen en una línea bajo el título:
%%   Fecha: ...  |  Duración: ...  |  Valor: ...
%%
\fecha{15 de abril de 2026}
\duracion{2 horas}
\valor{100 puntos}

\begin{document}

%% \imprimirtitulo
%%   Imprime el encabezado del documento: título y línea de metadatos
%%   (si alguno está definido). Debe ir inmediatamente después de
%%   \begin{document}.
%%
\imprimirtitulo

%% \datosestudiante
%%   Línea para nombre y carné del estudiante. Omitir en hojas de
%%   ejercicios que no requieren identificación individual.
%%
\datosestudiante

%% ── Entorno instrucciones ─────────────────────────────────────────────────────
%%
%% Caja enmarcada con lista numerada de indicaciones.
%% Cada indicación va con \item.
%%
%% Alternativa más simple (sin caja, texto libre):
%%   \begin{indicaciones}
%%     Texto de las indicaciones...
%%   \end{indicaciones}
%%
\begin{instrucciones}
  \item Este examen es individual, con libro y apuntes cerrados.
  \item Justifique cada respuesta. Una respuesta sin desarrollo no se
        calificará aunque sea correcta.
  \item Está prohibido el uso de calculadora.
  \item Si necesita más espacio, use el reverso de la hoja e indíquelo
        claramente.
  \item La deshonestidad académica se reportará según el Reglamento
        Estudiantil de la Universidad de Costa Rica.
\end{instrucciones}

%% ── Entorno ejercicio ────────────────────────────────────────────────────────
%%
%% \begin{ejercicio}[N]   N es el valor en puntos (opcional).
%% Numeración automática. Puede combinarse con \begin{subejercicios}
%% para incisos a), b), c), ...
%%
\begin{ejercicio}[25]
  Resuelva el siguiente sistema de ecuaciones aplicando eliminación de
  Gauss-Jordan. Dé el conjunto solución.
  \[
    \begin{cases}
       x +  y +  z =  6 \\
      2x + 3y +  z = 11 \\
       x + 2y + 3z = 14
    \end{cases}
  \]

  %% \begin{solucion}...\end{solucion}
  %%   Visible solo con la opción [soluciones]. Sin ella, no genera output.
  %%
  \begin{solucion}
    La matriz aumentada es
    \[
      \left(\begin{array}{ccc|r} 1 & 1 & 1 & 6 \\ 2 & 3 & 1 & 11 \\ 1 & 2 & 3 & 14
      \end{array}\right)
      \longrightarrow \cdots \longrightarrow
      \left(\begin{array}{ccc|r} 1 & 0 & 0 & 1 \\ 0 & 1 & 0 & 2 \\ 0 & 0 & 1 & 3
      \end{array}\right).
    \]
    El conjunto solución es $S = \{(1,\,2,\,3)\}$.
  \end{solucion}
\end{ejercicio}

\begin{ejercicio}[25]
  Determine el rango de cada matriz.

  \begin{subejercicios}
    \item $\begin{pmatrix} 1 & 2 & 3 \\ 2 & 5 & 7 \\ 1 & 3 & 4 \end{pmatrix}$

    \begin{solucion}
      El rango es $\mathbf{2}$.
    \end{solucion}

    \item $\begin{pmatrix} 1 & 3 & 2 \\ 2 & 6 & 4 \end{pmatrix}$

    \begin{solucion}
      El rango es $\mathbf{1}$.
    \end{solucion}
  \end{subejercicios}
\end{ejercicio}

\begin{ejercicio}[25]
  Determine si el siguiente sistema homogéneo tiene solución no trivial.
  Justifique.
  \[
    \begin{cases}
       x +  y +  z = 0 \\
      2x + 3y + 4z = 0 \\
       x + 2y + 3z = 0
    \end{cases}
  \]

  \begin{solucion}
    Reduciendo: el sistema tiene solución no trivial (variable libre).
    El conjunto solución es $S = \{(t,\,-2t,\,t) : t \in \mathbb{R}\}$.
  \end{solucion}
\end{ejercicio}

\begin{ejercicio}[25]
  Demuestre que si $A$ es una matriz de $3\times 3$ con rango~$3$, entonces el
  sistema $A\mathbf{x} = \mathbf{b}$ tiene solución única para todo
  $\mathbf{b} \in \mathbb{R}^3$.

  \begin{solucion}
    Si $\operatorname{rango}(A) = 3$, la forma escalonada reducida de $A$ es la
    identidad $I_3$. Para cualquier $\mathbf{b}$, la eliminación de Gauss-Jordan
    produce una única solución $\mathbf{x} = A^{-1}\mathbf{b}$.
  \end{solucion}
\end{ejercicio}

\end{document}

%%% Local Variables:
%%% mode: LaTeX
%%% TeX-master: t
%%% End:

%%
\documentclass{emate-ucr}
% \documentclass[soluciones]{emate-ucr}   % <-- versión con soluciones

%% ── Metadatos obligatorios ───────────────────────────────────────────────────
%%
%% \curso{texto}    Centro del encabezado de página. \\ para saltos de línea.
%% \encabezado{texto}  Título en negrita al inicio del documento.
%%
\curso{I Ciclo 2026\\MA-1022\\Cálculo para Ciencias Económicas II}
\encabezado{Primer Examen Parcial}

%% ── Metadatos opcionales ─────────────────────────────────────────────────────
%%
%% Cualquier subconjunto de estos tres puede usarse o no.
%% Los que estén definidos aparecen en una línea bajo el título:
%%   Fecha: ...  |  Duración: ...  |  Valor: ...
%%
\fecha{15 de abril de 2026}
\duracion{2 horas}
\valor{100 puntos}

\begin{document}

%% \imprimirtitulo
%%   Imprime el encabezado del documento: título y línea de metadatos
%%   (si alguno está definido). Debe ir inmediatamente después de
%%   \begin{document}.
%%
\imprimirtitulo

%% \datosestudiante
%%   Línea para nombre y carné del estudiante. Omitir en hojas de
%%   ejercicios que no requieren identificación individual.
%%
\datosestudiante

%% ── Entorno instrucciones ─────────────────────────────────────────────────────
%%
%% Caja enmarcada con lista numerada de indicaciones.
%% Cada indicación va con \item.
%%
%% Alternativa más simple (sin caja, texto libre):
%%   \begin{indicaciones}
%%     Texto de las indicaciones...
%%   \end{indicaciones}
%%
\begin{instrucciones}
  \item Este examen es individual, con libro y apuntes cerrados.
  \item Justifique cada respuesta. Una respuesta sin desarrollo no se
        calificará aunque sea correcta.
  \item Está prohibido el uso de calculadora.
  \item Si necesita más espacio, use el reverso de la hoja e indíquelo
        claramente.
  \item La deshonestidad académica se reportará según el Reglamento
        Estudiantil de la Universidad de Costa Rica.
\end{instrucciones}

%% ── Entorno ejercicio ────────────────────────────────────────────────────────
%%
%% \begin{ejercicio}[N]   N es el valor en puntos (opcional).
%% Numeración automática. Puede combinarse con \begin{subejercicios}
%% para incisos a), b), c), ...
%%
\begin{ejercicio}[25]
  Resuelva el siguiente sistema de ecuaciones aplicando eliminación de
  Gauss-Jordan. Dé el conjunto solución.
  \[
    \begin{cases}
       x +  y +  z =  6 \\
      2x + 3y +  z = 11 \\
       x + 2y + 3z = 14
    \end{cases}
  \]

  %% \begin{solucion}...\end{solucion}
  %%   Visible solo con la opción [soluciones]. Sin ella, no genera output.
  %%
  \begin{solucion}
    La matriz aumentada es
    \[
      \left(\begin{array}{ccc|r} 1 & 1 & 1 & 6 \\ 2 & 3 & 1 & 11 \\ 1 & 2 & 3 & 14
      \end{array}\right)
      \longrightarrow \cdots \longrightarrow
      \left(\begin{array}{ccc|r} 1 & 0 & 0 & 1 \\ 0 & 1 & 0 & 2 \\ 0 & 0 & 1 & 3
      \end{array}\right).
    \]
    El conjunto solución es $S = \{(1,\,2,\,3)\}$.
  \end{solucion}
\end{ejercicio}

\begin{ejercicio}[25]
  Determine el rango de cada matriz.

  \begin{subejercicios}
    \item $\begin{pmatrix} 1 & 2 & 3 \\ 2 & 5 & 7 \\ 1 & 3 & 4 \end{pmatrix}$

    \begin{solucion}
      El rango es $\mathbf{2}$.
    \end{solucion}

    \item $\begin{pmatrix} 1 & 3 & 2 \\ 2 & 6 & 4 \end{pmatrix}$

    \begin{solucion}
      El rango es $\mathbf{1}$.
    \end{solucion}
  \end{subejercicios}
\end{ejercicio}

\begin{ejercicio}[25]
  Determine si el siguiente sistema homogéneo tiene solución no trivial.
  Justifique.
  \[
    \begin{cases}
       x +  y +  z = 0 \\
      2x + 3y + 4z = 0 \\
       x + 2y + 3z = 0
    \end{cases}
  \]

  \begin{solucion}
    Reduciendo: el sistema tiene solución no trivial (variable libre).
    El conjunto solución es $S = \{(t,\,-2t,\,t) : t \in \mathbb{R}\}$.
  \end{solucion}
\end{ejercicio}

\begin{ejercicio}[25]
  Demuestre que si $A$ es una matriz de $3\times 3$ con rango~$3$, entonces el
  sistema $A\mathbf{x} = \mathbf{b}$ tiene solución única para todo
  $\mathbf{b} \in \mathbb{R}^3$.

  \begin{solucion}
    Si $\operatorname{rango}(A) = 3$, la forma escalonada reducida de $A$ es la
    identidad $I_3$. Para cualquier $\mathbf{b}$, la eliminación de Gauss-Jordan
    produce una única solución $\mathbf{x} = A^{-1}\mathbf{b}$.
  \end{solucion}
\end{ejercicio}

\end{document}

%%% Local Variables:
%%% mode: LaTeX
%%% TeX-master: t
%%% End:

%%
\documentclass{emate-ucr}
% \documentclass[soluciones]{emate-ucr}   % <-- versión con soluciones

%% ── Metadatos obligatorios ───────────────────────────────────────────────────
%%
%% \curso{texto}    Centro del encabezado de página. \\ para saltos de línea.
%% \encabezado{texto}  Título en negrita al inicio del documento.
%%
\curso{I Ciclo 2026\\MA-1022\\Cálculo para Ciencias Económicas II}
\encabezado{Primer Examen Parcial}

%% ── Metadatos opcionales ─────────────────────────────────────────────────────
%%
%% Cualquier subconjunto de estos tres puede usarse o no.
%% Los que estén definidos aparecen en una línea bajo el título:
%%   Fecha: ...  |  Duración: ...  |  Valor: ...
%%
\fecha{15 de abril de 2026}
\duracion{2 horas}
\valor{100 puntos}

\begin{document}

%% \imprimirtitulo
%%   Imprime el encabezado del documento: título y línea de metadatos
%%   (si alguno está definido). Debe ir inmediatamente después de
%%   \begin{document}.
%%
\imprimirtitulo

%% \datosestudiante
%%   Línea para nombre y carné del estudiante. Omitir en hojas de
%%   ejercicios que no requieren identificación individual.
%%
\datosestudiante

%% ── Entorno instrucciones ─────────────────────────────────────────────────────
%%
%% Caja enmarcada con lista numerada de indicaciones.
%% Cada indicación va con \item.
%%
%% Alternativa más simple (sin caja, texto libre):
%%   \begin{indicaciones}
%%     Texto de las indicaciones...
%%   \end{indicaciones}
%%
\begin{instrucciones}
  \item Este examen es individual, con libro y apuntes cerrados.
  \item Justifique cada respuesta. Una respuesta sin desarrollo no se
        calificará aunque sea correcta.
  \item Está prohibido el uso de calculadora.
  \item Si necesita más espacio, use el reverso de la hoja e indíquelo
        claramente.
  \item La deshonestidad académica se reportará según el Reglamento
        Estudiantil de la Universidad de Costa Rica.
\end{instrucciones}

%% ── Entorno ejercicio ────────────────────────────────────────────────────────
%%
%% \begin{ejercicio}[N]   N es el valor en puntos (opcional).
%% Numeración automática. Puede combinarse con \begin{subejercicios}
%% para incisos a), b), c), ...
%%
\begin{ejercicio}[25]
  Resuelva el siguiente sistema de ecuaciones aplicando eliminación de
  Gauss-Jordan. Dé el conjunto solución.
  \[
    \begin{cases}
       x +  y +  z =  6 \\
      2x + 3y +  z = 11 \\
       x + 2y + 3z = 14
    \end{cases}
  \]

  %% \begin{solucion}...\end{solucion}
  %%   Visible solo con la opción [soluciones]. Sin ella, no genera output.
  %%
  \begin{solucion}
    La matriz aumentada es
    \[
      \left(\begin{array}{ccc|r} 1 & 1 & 1 & 6 \\ 2 & 3 & 1 & 11 \\ 1 & 2 & 3 & 14
      \end{array}\right)
      \longrightarrow \cdots \longrightarrow
      \left(\begin{array}{ccc|r} 1 & 0 & 0 & 1 \\ 0 & 1 & 0 & 2 \\ 0 & 0 & 1 & 3
      \end{array}\right).
    \]
    El conjunto solución es $S = \{(1,\,2,\,3)\}$.
  \end{solucion}
\end{ejercicio}

\begin{ejercicio}[25]
  Determine el rango de cada matriz.

  \begin{subejercicios}
    \item $\begin{pmatrix} 1 & 2 & 3 \\ 2 & 5 & 7 \\ 1 & 3 & 4 \end{pmatrix}$

    \begin{solucion}
      El rango es $\mathbf{2}$.
    \end{solucion}

    \item $\begin{pmatrix} 1 & 3 & 2 \\ 2 & 6 & 4 \end{pmatrix}$

    \begin{solucion}
      El rango es $\mathbf{1}$.
    \end{solucion}
  \end{subejercicios}
\end{ejercicio}

\begin{ejercicio}[25]
  Determine si el siguiente sistema homogéneo tiene solución no trivial.
  Justifique.
  \[
    \begin{cases}
       x +  y +  z = 0 \\
      2x + 3y + 4z = 0 \\
       x + 2y + 3z = 0
    \end{cases}
  \]

  \begin{solucion}
    Reduciendo: el sistema tiene solución no trivial (variable libre).
    El conjunto solución es $S = \{(t,\,-2t,\,t) : t \in \mathbb{R}\}$.
  \end{solucion}
\end{ejercicio}

\begin{ejercicio}[25]
  Demuestre que si $A$ es una matriz de $3\times 3$ con rango~$3$, entonces el
  sistema $A\mathbf{x} = \mathbf{b}$ tiene solución única para todo
  $\mathbf{b} \in \mathbb{R}^3$.

  \begin{solucion}
    Si $\operatorname{rango}(A) = 3$, la forma escalonada reducida de $A$ es la
    identidad $I_3$. Para cualquier $\mathbf{b}$, la eliminación de Gauss-Jordan
    produce una única solución $\mathbf{x} = A^{-1}\mathbf{b}$.
  \end{solucion}
\end{ejercicio}

\end{document}

%%% Local Variables:
%%% mode: LaTeX
%%% TeX-master: t
%%% End:
